% Methodology Section for MobilePIXORNeXt
% Written according to Springer Nature - Arabian Journal for Science and Engineering (AJSE)
% Structure: Module Triad (Motivation -> Design -> Technical Advantage) + Equation Closure + Pseudocode Algorithm

\begin{strip}
    \centering
    \paperfigure[width=0.88\linewidth]{pipeline_overview.pdf}{32mm}{Overall architecture and data flow of MobilePIXORNeXt. A raw LiDAR point cloud is converted via the parameter-free Strata-BEV14 rasterizer into a 14-channel statistical tensor. The compact ConvNeXt-style backbone with Block-SimAM extracts multi-scale spatial features, which are fused by the Scale-Gated Feature Pyramid Network (SG-FPN). Decoupled grouped prediction towers output heatmaps, subcell offsets, log-sizes, doubled yaw angles, and localization quality. In training, the network is supervised by Oriented Geometric Alignment (OGA) and IoU-Aware Quality Alignment (IQA), dynamically balanced via Temperature-Softmax Bounded Uncertainty Weighting (T-SBUW).}
    \captionof{figure}{End-to-end architecture and tensor flow of MobilePIXORNeXt. Unordered 3D LiDAR point clouds pass through the parameter-free Strata-BEV14 rasterizer into a 14-channel statistical BEV tensor $\mathbf{X} \in \mathbb{R}^{B \times 14 \times 800 \times 704}$ without learned pillar encoders. ConvNeXt-style stages with parameter-free Block-SimAM attention extract multi-scale features. A top-down Scale-Gated Feature Pyramid Network (SG-FPN) fuses them into a 32-channel shared representation at stride 4 ($200 \times 176$). Decoupled GroupedHeader towers process Car and Pedestrian/Cyclist categories independently. Solid arrows show the feedforward inference path. Dashed arrows show training-only Oriented Geometric Alignment (OGA), IoU-Aware Quality Alignment (IQA), and Temperature-Softmax Bounded Uncertainty Weighting (T-SBUW).}
    \label{fig:r37_overview}
\end{strip}

MobilePIXORNeXt is a compact, proposal-free, multi-class three-dimensional (3D) object detector for real-time Light Detection and Ranging (LiDAR) perception on resource-constrained embedded platforms. Its Bird's-Eye View (BEV) representation avoids learned point-cloud encoders and three-dimensional sparse convolutions. Figure~\ref{fig:r37_overview} shows the data flow from the raw point cloud to decoded oriented bounding boxes.

\par\vspace{-0.5\baselineskip}
\subsection{Problem Formulation and Input/Output Contracts}
\label{sec:problem_formulation}

\paragraph{LiDAR Coordinate Frame and Spatial Region of Interest.}
An unordered LiDAR point cloud acquired by a rotating scanner is represented as $\mathcal{P} = \{\mathbf{p}_i\}_{i=1}^{N_{\mathrm{pts}}}$. Each point $\mathbf{p}_i = (x_i, y_i, z_i, \iota_i) \in \mathbb{R}^4$ contains longitudinal distance $x_i$ (forward positive), lateral distance $y_i$ (leftward positive), vertical elevation $z_i$ (upward positive), and calibrated reflectance intensity $\iota_i \in [0, 1]$. Perception is restricted to a Cartesian Region of Interest (ROI):
\begin{equation}
    \begin{aligned}
        x_i &\in [x_{\min}, x_{\max}], \\
        y_i &\in [y_{\min}, y_{\max}], \\  z_i &\in [z_{\min}, z_{\max}].
    \end{aligned}
    \label{eq:spatial_roi}
\end{equation}
We used longitudinal bounds $[0.0, 70.4]$\,m, lateral bounds $[-40.0, 40.0]$\,m, and vertical bounds $[-2.5, 1.0]$\,m. An open-boundary inward margin of $\epsilon = 0.001$\,m reduced boundary artifacts. A horizontal resolution of $\Delta_x = \Delta_y = 0.1$\,m gave a grid of $H_{\mathrm{in}} \times W_{\mathrm{in}} = 800 \times 704$ cells. A point $\mathbf{p}_i$ maps to grid coordinate $(u_i, v_i)$ as:
\begin{equation}
    u_i = \left\lfloor \frac{x_i - x_{\min}}{\Delta_x} \right\rfloor, \qquad v_i = \left\lfloor \frac{y_i - y_{\min}}{\Delta_y} \right\rfloor.
    \label{eq:grid_assignment}
\end{equation}

\paragraph{Detection Target Contract.}
The detector predicts oriented bounding boxes projected onto the horizontal BEV plane across three semantic categories: Car ($c=0$), Pedestrian ($c=1$), and Cyclist ($c=2$). At an output spatial downsampling factor of $s_{\mathrm{out}} = 4$, the output feature map resolution is $H_{\mathrm{out}} \times W_{\mathrm{out}} = 200 \times 176$. For each detected object instance $k$, the target state is parameterized by the 7-tuple:
\begin{equation}
    \mathbf{b}_k = (c_k, x_k^*, y_k^*, l_k^*, w_k^*, \psi_k^*, s_k),
    \label{eq:target_tuple}
\end{equation}
where $(x_k^*, y_k^*) \in \mathbb{R}^2$ denotes continuous metric center coordinates and $(l_k^*, w_k^*) \in \mathbb{R}^2$ denotes length and width. The heading $\psi_k^* \in [-\pi/2, \pi/2]$ uses a $\pi$-symmetric doubled-angle representation, and $s_k \in [0, 1]$ is the calibrated confidence score. These targets require a dense BEV representation that retains fine 3D geometry without a learned point encoder.

\par\vspace{-0.5\baselineskip}
\subsection{Deterministic Strata-BEV14 Feature Representation}
\label{sec:strata_bev14}

\par\vspace{-0.5\baselineskip}
\noindent
\begin{minipage}{\columnwidth}
    \centering
    \paperfigure{strata_bev14.pdf}{10mm}{Visual representation of the fourteen deterministic feature channels in Strata-BEV14 for a representative LiDAR scene. Channels 0–3 capture stratified height-band densities, Channels 4–7 encode vertical elevation statistics (peak, mean, span, deviation), Channels 8–9 capture material surface reflectance, Channels 10–11 represent global density and occupancy, and Channels 12–13 encode continuous subcell horizontal offsets to eliminate grid discretization error.}
    \captionof{figure}{Fourteen deterministic feature channels in the Strata-BEV14 representation for a KITTI driving scene ($800 \times 704$ grid at $0.1\,\text{m}$ resolution). Channels 0--3 encode logarithmic point densities across four vertical strata ($[-2.5, 1.0]\,\text{m}$). Channels 4--7 encode normalized maximum elevation $z_{\max}$, mean elevation $z_{\mathrm{mean}}$, elevation span $z_{\mathrm{span}}$, and population standard deviation $z_{\mathrm{std}}$. Channels 8--9 encode maximum and mean reflectance intensity. Channel 10 records total cell logarithmic density, and Channel 11 records binary occupancy. Channels 12--13 preserve continuous horizontal subcell centroid displacements $(\Delta\bar{x}, \Delta\bar{y}) \in [-0.5, 0.5]$ to reduce discretization error. Brighter values indicate higher point concentrations, reflectance, or vertical dispersion.}
    \label{fig:strata_bev14}
\end{minipage}
\par

\paragraph{Motivation.}
Voxel- and pillar-based detectors~\cite{lang2019pointpillars,li2023pillarnext} use PointNet-style column encoders, which add computation and padding overhead. 35-channel binary slicing~\cite{yang2018pixor} omits continuous height and reflectance cues, while three-channel height bands~\cite{khoshkdahan2026triband} omit higher-order geometry. We developed \textbf{Strata-BEV14} (Figure~\ref{fig:strata_bev14}), a deterministic, parameter-free rasterizer, to retain these features with low computational cost.


\paragraph{Mathematical Formulation.}
For each occupied cell $(u, v)$ with $N$ points $\mathcal{P}_{u, v} = \{\mathbf{p}_i\}_{i=1}^N$, we normalized vertical elevations to $\zeta_i = (z_i - z_{\min}) / (z_{\max} - z_{\min}) \in [0, 1]$ and clamped reflectance values $\iota_i$ to $[0, 1]$. Strata-BEV14 computes fourteen analytical features in four groups (Table~\ref{tab:strata_bev14}):

1. \textit{Height-Stratified Density Channels ($F_0$--$F_3$)}: Normalized height $\zeta_i$ is partitioned into four uniform vertical intervals $b \in \{0, 1, 2, 3\}$. For bin point count $N_b$, a logarithmic compression is applied to mitigate LiDAR ray density fall-off over range:
\begin{equation}
    F_b = d_{\log}(N_b) = \min\left(1, \frac{\ln(1 + N_b)}{\ln(1 + N_{\mathrm{norm}})}\right),
    \label{eq:density_encoding}
\end{equation}
where $N_{\mathrm{norm}} = 32$ is a normalization constant chosen to saturate at dense returns.

2. \textit{Vertical Spatial Moments ($F_4$--$F_7$)}: To encode obstacle geometry without explicit voxelization, we extract peak elevation $F_4 = \max_i \zeta_i$, mean elevation $F_5 = \frac{1}{N}\sum_i \zeta_i$, vertical span $F_6 = \max_i \zeta_i - \min_i \zeta_i$, and the population standard deviation:
\begin{equation}
    F_7 = \begin{cases}
        \sqrt{\max\left(0, \frac{1}{N} \sum_{i=1}^N \zeta_i^2 - F_5^2\right)}, & \text{if } N > 1, \\
        0, & \text{otherwise}.
    \end{cases}
    \label{eq:z_std_encoding}
\end{equation}

3. \textit{Reflectance Statistics ($F_8$--$F_9$)}: Calibrated surface returns capture material discrimination (e.g., distinguishing retro-reflective license plates, pedestrian clothing, or road signs):
\begin{equation}
    F_8 = \max_{i=1,\dots,N} \iota_i, \qquad F_9 = \frac{1}{N} \sum_{i=1}^N \iota_i.
    \label{eq:intensity_encoding}
\end{equation}

4. \textit{Occupancy and Continuous Subcell Offsets ($F_{10}$--$F_{13}$)}: Channel $F_{10} = d_{\log}(N)$ records total point density, and $F_{11} = 1.0$ indicates occupancy. For $0.1$\,m cells, channels $F_{12}$ and $F_{13}$ retain continuous horizontal subcell centroid displacements to reduce quantization error:
\begin{align}
    x_u^{\mathrm{cell}} &= x_{\min} + (u + 0.5)\Delta_x, \nonumber\\
    y_v^{\mathrm{cell}} &= y_{\min} + (v + 0.5)\Delta_y, \nonumber\\
    F_{12} &= \operatorname{clip}\left(\frac{\sum_{i=1}^N (x_i - x_u^{\mathrm{cell}})}{N\Delta_x}, -0.5, 0.5\right), \nonumber\\
    F_{13} &= \operatorname{clip}\left(\frac{\sum_{i=1}^N (y_i - y_v^{\mathrm{cell}})}{N\Delta_y}, -0.5, 0.5\right).
    \label{eq:offset_encoding}
\end{align}
Empty cells are assigned zeros across all 14 channels, producing the structured input tensor $\mathbf{X} \in \mathbb{R}^{B \times 14 \times 800 \times 704}$.

\begin{table}[t]
    \caption{Mathematical formulations and physical semantics of the fourteen deterministic feature channels in the Strata-BEV14 representation. All statistical moments and spatial offsets are evaluated deterministically over LiDAR points within horizontal grid cells of size $\Delta_x = \Delta_y = 0.1\,\text{m}$ ($800 \times 704$ grid) using NumPy or Numba on CPU without learned parameters.}
    \label{tab:strata_bev14}
    \centering\small
    \setlength{\tabcolsep}{3pt}
    \begin{tabular}{@{}c >{\raggedright\arraybackslash}p{0.20\columnwidth} >{\raggedright\arraybackslash}p{0.55\columnwidth}@{}}
        \toprule
        \textbf{Channel} & \textbf{Feature Identity} & \textbf{Mathematical Definition and Semantics} \\
        \midrule
        $0$--$3$ & Stratified Log Counts & Four height-band densities: $F_b = d_{\log}(N_b)$, where $b = \min(\lfloor 4\zeta_i \rfloor, 3)$. \\
        $4$ & Maximum Elevation & $F_4 = \max_{i} \zeta_i$, peak normalized elevation. \\
        $5$ & Mean Elevation & $F_5 = \frac{1}{N} \sum_{i=1}^N \zeta_i$, first vertical moment. \\
        $6$ & Elevation Span & $F_6 = \max_{i} \zeta_i - \min_{i} \zeta_i$, continuous vertical extent. \\
        $7$ & Elevation Deviation & $F_7 = \sqrt{\max(0, \frac{1}{N} \sum_{i} \zeta_i^2 - F_5^2)}$, spatial dispersion. \\
        $8$ & Peak Reflectance & $F_8 = \max_{i} \iota_i$, highest material reflectivity. \\
        $9$ & Mean Reflectance & $F_9 = \frac{1}{N} \sum_{i=1}^N \iota_i$, average surface intensity. \\
        $10$ & Global Density & $F_{10} = d_{\log}(N)$, overall point density. \\
        $11$ & Binary Occupancy & $F_{11} = 1.0$ if $N \ge 1$, and $0.0$ otherwise. \\
        $12$--$13$ & Centroid Displacements & Subcell offsets $\Delta \bar{x}, \Delta \bar{y} \in [-0.5, 0.5]$ relative to cell center. \\
        \bottomrule
    \end{tabular}
\end{table}

\paragraph{Technical Advantage.}
Strata-BEV14 provides both a NumPy reference implementation and a Numba-accelerated CPU rasterizer. Both compute identical fourteen deterministic feature channels without learned parameters. The Numba backend accumulates point statistics, normalizes occupied cells, and formats the output into CHW layout for direct backbone ingestion. In the measured single-threaded pipeline, point-cloud loading and rasterization together require $20.94\,\text{ms}$, encompassing disk I/O, coordinate parsing, and cell aggregation. Because point-to-cell binning and statistic accumulation (extents, mean intensity, log-density lookups) operate on spatially disjoint grid cells with zero inter-cell data dependencies, the rasterization algorithm is embarrassingly parallel. On embedded System-on-Chip (SoC) architectures (e.g., NVIDIA Jetson AGX Orin), this deterministic accumulation can be mapped directly to an optimized CUDA kernel executing in under $3\,\text{ms}$. Alternatively, in vehicular streaming deployment, CPU rasterization operates via asynchronous double-buffering: a companion CPU core rasterizes scan $t+1$ concurrently with GPU neural network inference on scan $t$, effectively hiding preprocessing latency behind the GPU execution window.

\subsection{Physics-Consistent Ground-Truth Augmentation}
\label{sec:pc_gta}

\paragraph{Motivation.}
Cut-and-paste ground-truth (GT) augmentation balances rare LiDAR categories, but naive insertion ignores scene physics. It can produce floating or submerged objects, background collisions, missing raycast shadows, and unrealistic range-dependent point densities. We developed \textbf{Physics-Consistent Ground-Truth Augmentation} (\texttt{PC-GTA}) to address these artifacts. It improved Moderate mAP by $+6.03\,\text{pp}$ over naive sampling in the ablations.

\paragraph{Design Heuristics.}
The sampler provides five consistency mechanisms when placing a box $\mathbf{b}_{\mathrm{insert}} = (x, y, z, l, w, h, \psi)$ at range $d = \sqrt{x^2 + y^2}$. The mechanisms can be enabled independently; the settings used for the reported final run are given in Section~\ref{sec:implementation_details}:
\begin{enumerate}
\item \textit{Adaptive Ground Plane Fitting}: To prevent floating or sunken objects, local road elevation is estimated within a range-dependent search radius:
\begin{align}
    R(d) &= \operatorname{clip}(2.5 + 0.03\max(0, d - 15), 2.5, 4.0)\,\mathrm{m},\nonumber\\[0.8em]
    \hat{z}_{\mathrm{ground}} &= \operatorname{percentile}(\{z_j \mid \mathbf{p}_j \in \mathcal{N}(x, y)\}, 5),
    \label{eq:ground_fitting}
\end{align}
where $\mathcal{N}(x, y) = \{\mathbf{p}_j \mid (x_j - x)^2 + (y_j - y)^2 \le R(d)^2\}$ denotes the local neighborhood used for ground estimation.
The inserted object elevation is anchored flush with the estimated local ground elevation: $z_{\mathrm{insert}} = \hat{z}_{\mathrm{ground}} + 0.5h$.

\item \textit{Static Obstacle Clearance}: To prevent clipping into physical structures, a candidate box is rejected if existing background points within its bounding footprint exceed a strict clearance threshold: $N_{\mathrm{static}} > N_{\mathrm{thresh}}$ ($N_{\mathrm{thresh}} = 2$; valid placement requires $N_{\mathrm{static}} \le N_{\mathrm{thresh}}$).

\item \textit{Optional Line-of-Sight (LOS) Occlusion Check}: When enabled, the optical ray from sensor origin to box center $(x, y)$ is cast. If existing dense clusters obstruct the trajectory, the insertion candidate is discarded.

\item \textit{Optional Raycast LiDAR Shadow Masking}: In LiDAR scans, foreground objects occlude returns from objects behind them. When enabled, this mechanism computes intersections for each inserted box using rays $\mathbf{r}(t) = t \cdot \mathbf{p}^{xyz}$ through background points with 3D metric coordinates $\mathbf{p}^{xyz} = (x_i, y_i, z_i) \in \mathbb{R}^3$, decoupling reflectance. It masks points behind the exit face ($t_{\mathrm{leave}} \le 1.0$) whose radial distance satisfies $(1 - t_{\mathrm{leave}})\|\mathbf{p}^{xyz}\| > 0.05$\,m:
\begin{equation}
    \begin{aligned}
        \mathcal{P}_{\mathrm{augmented}}
        &= \mathcal{P}_{\mathrm{target}} \setminus
        \{\mathbf{p} \in \mathcal{P}_{\mathrm{target}} \mid \\
        &\qquad \operatorname{Mask}_{\mathrm{shadow}}(\mathbf{p}) = \mathrm{True}\}
        \cup \mathcal{P}_{\mathrm{object}}.
    \end{aligned}
    \label{eq:shadow_mask}
\end{equation}

\item \textit{Range-Adaptive Point and Intensity Attenuation}: If an object captured at source distance $r_{\mathrm{src}}$ is placed at target distance $r_{\mathrm{tgt}} > r_{\mathrm{src}}$, its point cloud is randomly subsampled according to the inverse-square divergence law:
\begin{equation}
    P_{\mathrm{keep}} = \min\left(1, \left(\frac{r_{\mathrm{src}}}{r_{\mathrm{tgt}}}\right)^2\right),
    \label{eq:inverse_square}
\end{equation}
and the optional radiometric adjustment attenuates reflectance intensity according to:
\[
    \begin{aligned}
        \eta_{\mathrm{int}}
        &= \left(\frac{r_{\mathrm{src}}}{r_{\mathrm{tgt}}}\right)^{\max(0, 2 - \gamma)},
        \qquad \gamma = 1.7, \\
        I_{\mathrm{tgt}}
        &= \operatorname{clip}(I_{\mathrm{src}}\eta_{\mathrm{int}}, 0, 1).
    \end{aligned}
\]
\end{enumerate}
During training, the resulting physically consistent point cloud $\mathcal{P}_{\mathrm{augmented}}$ is rasterized into the 14-channel input tensor $\mathbf{X}$ and fed into the convolutional backbone hierarchy.

\subsection{ConvNeXt-Style Multi-Scale BEV Backbone Hierarchy}
\label{sec:backbone_architecture}

\begin{figure*}[t]
    \centering
    \paperfigure[width=0.95\textwidth]{backbone_f32.pdf}{42mm}{Detailed structural topology of the MobilePIXORNeXt feedforward architecture, comprising the multi-scale backbone hierarchy, Top-Down Scale-Gated Neck (SG-FPN), and decoupled grouped prediction towers. Spatial feature maps are progressively downsampled at strides 2, 4, 8, and 16 using repeated MobilePIXORNeXt-S blocks. Scale-gated top-down fusion balances deep semantics with fine spatial localization into a shared 32-channel representation at stride 4 ($200 \times 176$). Independent grouped prediction towers decouple Car from Pedestrian and Cyclist.}
    \caption{MobilePIXORNeXt feedforward hierarchy with layer dimensions, channel widths, and receptive fields. A two-layer stem processes 14-channel input at stride 2 ($400 \times 352$). Three stages follow: C2 has 48 channels and depth 3; C3 has 96 channels and depth 4; C4 has 128 channels and depth 2. Each stage uses \texttt{MobilePIXORNeXt-S} blocks with Block-SimAM. The Scale-Gated FPN (SG-FPN) uses lateral $1 \times 1$ projections, bilinear $2\times$ upsampling, depthwise $3 \times 3$ refinement, and learnable scale gates $\mathbf{G}_r$. It produces a shared 32-channel feature map ($F_{32}$) at stride 4 ($200 \times 176$). Independent grouped towers separate Car predictions from Pedestrian and Cyclist predictions to reduce aspect-ratio conflicts within the parameter budget (753,569 total parameters).}
    \label{fig:backbone}
\end{figure*}


\paragraph{Motivation.}
BEV object footprints range from pedestrians ($l \approx 0.6$\,m) and cyclists ($l \approx 1.8$\,m) to passenger vehicles ($l \approx 4.5$--$5.0$\,m). Small $3 \times 3$ convolutions require deeper networks to cover this spatial context, increasing latency and parameter count. At an input resolution of $\Delta_x = \Delta_y = 0.1$\,m, a $7 \times 7$ depthwise receptive field spans $0.7 \times 0.7$\,m at stride 1. It expands to over $5.6 \times 5.6$\,m by Stage 3 (stride 8), enough to cover a passenger vehicle without dilation. We designed \textbf{MobilePIXORNeXt Backbone} with large $7 \times 7$ depthwise convolutions, inverted bottleneck expansions, and LayerScale stabilization in repeated \texttt{MobilePIXORNeXt-S} blocks to obtain this context within a small parameter budget.

\paragraph{Backbone Topology and Multi-Stage Hierarchy.}
As depicted in Figure~\ref{fig:backbone}, the network topology comprises:
\begin{enumerate}
    \item \textit{Convolutional Stem}: Input $\mathbf{X} \in \mathbb{R}^{B \times 14 \times 800 \times 704}$ passes through two consecutive $3 \times 3$ convolutions, with stride 2 followed by stride 1 and 32 output channels. Batch Normalization (BN) and Sigmoid Linear Unit (SiLU) activations produce stride-2 representation $\mathbf{F}^{(2)} \in \mathbb{R}^{B \times 32 \times 400 \times 352}$.
    \item \textit{Progressive Multi-Stage Hierarchy}: Three progressive stages are separated by strided transition downsample blocks (\texttt{DownsampleBlock}: Conv $3 \times 3$ stride 2 + BN + SiLU):
    \begin{itemize}
        \item Stage 2 (Stride 4): 3 \texttt{MobilePIXORNeXt-S} blocks, $C_2 = 48 \rightarrow \mathbf{F}^{(4)} \in \mathbb{R}^{B \times 48 \times 200 \times 176}$.
        \item Stage 3 (Stride 8): 4 \texttt{MobilePIXORNeXt-S} blocks, $C_3 = 96 \rightarrow \mathbf{F}^{(8)} \in \mathbb{R}^{B \times 96 \times 100 \times 88}$.
        \item Stage 4 (Stride 16): 2 \texttt{MobilePIXORNeXt-S} blocks, $C_4 = 128 \rightarrow \mathbf{F}^{(16)} \in \mathbb{R}^{B \times 128 \times 50 \times 44}$.
    \end{itemize}
\end{enumerate}

\paragraph{Streamlined Single-Branch Depthwise Design.}
MobilePIXORNeXt uses a single-branch large-kernel ($7 \times 7$) depthwise convolution during both training and deployment. It avoids the multi-path training topology and post-hoc weight folding of multi-branch reparameterization. Inverted bottleneck expansion ($e=2.5$) and LayerScale stabilization support optimization without auxiliary-branch memory during backpropagation. The backbone achieved $55.94$\,FPS on Tesla T4 and contained $604,864$ parameters.

\subsection{MobilePIXORNeXt-S Residual Block with Embedded Block-SimAM}
\label{sec:next_s_block}

\begin{figure*}[t]
    \centering
    \paperfigure[width=0.95\linewidth,height=58mm,keepaspectratio]{next_simam_block.pdf}{58mm}{Micro-architecture of the MobilePIXORNeXt-S residual block and embedded parameter-free Block-SimAM energy attention. (Top) Complete feedforward path from input X to output Y through the 7x7 depthwise convolution, inverted bottleneck expansion (e=2.5), SimAM modulation, LayerScale scaling, and identity residual addition. (Bottom) Zoom-in structure of the parameter-free SimAM module, detailing spatial moment estimation, inverse-energy importance calculation, sigmoid saliency mapping, and element-wise feature recalibration.}
    \caption{The \texttt{MobilePIXORNeXt-S} residual block and parameter-free Block-SimAM attention. \textbf{Top}: Residual path from input $\mathbf{Z}$ to output $\mathbf{Z}'$. It applies a $7 \times 7$ depthwise convolution, inverted bottleneck expansion with ratio $e = 2.5$, and SiLU activation, followed by pointwise projection. Parameter-free 3D spatial-energy attention $\mathcal{S}(\mathbf{R})$ precedes LayerScale scaling $\boldsymbol{\gamma}_{\mathrm{blk}}$ (initialized to $10^{-5}$) and the identity shortcut. \textbf{Bottom}: SimAM computes channel-wise spatial moments ($\mu_c, \nu_c$) in FP32 precision and inverse-energy importance $s_t = (t - \mu_c)^2 / (4(\nu_c + \lambda_{\mathrm{sim}})) + 0.5$. Sigmoid scaling produces the 3D modulation tensor $\mathbf{A} \in [0, 1]^{B \times C \times H \times W}$ without learned parameters.}
    \label{fig:next_s_block}
\end{figure*}

\paragraph{Micro-Architecture and Inverted Bottleneck.}
Figure~\ref{fig:next_s_block} illustrates the \texttt{MobilePIXORNeXt-S} residual block. Given input tensor $\mathbf{Z} \in \mathbb{R}^{B \times C \times H \times W}$, each block applies an inverted bottleneck with expansion ratio $e = 2.5$:
\begin{align}
    \mathbf{V} &= \operatorname{BN}(\operatorname{DWConv}_{7\times7}(\mathbf{Z})),\nonumber\\
    \mathbf{R} &= \operatorname{BN}\left(W_{\mathrm{red}} \operatorname{SiLU}(W_{\mathrm{exp}}\mathbf{V})\right),\nonumber\\
    \mathbf{Z}' &= \mathbf{Z} + \boldsymbol{\gamma}_{\mathrm{blk}} \odot \mathcal{S}(\mathbf{R}).
    \label{eq:next_block}
\end{align}
Here, $W_{\mathrm{exp}} \in \mathbb{R}^{2.5C \times C \times 1 \times 1}$ and $W_{\mathrm{red}} \in \mathbb{R}^{C \times 2.5C \times 1 \times 1}$ are pointwise $1 \times 1$ convolutions. The operator $\mathcal{S}(\cdot)$ represents our Block-SimAM attention module, while $\boldsymbol{\gamma}_{\mathrm{blk}} \in \mathbb{R}^C$ is a learned LayerScale parameter initialized to $10^{-5}$ to ensure numerical stability during deep backpropagation.

\paragraph{Parameter-Free Block-SimAM Spatial-Energy Attention.}
Efficient Channel Attention (ECA-Net)~\cite{wang2020eca} and Squeeze-and-Excitation compute 1D channel weights by spatial pooling, applying the same weight to every location in a channel. Non-local spatial attention instead has quadratic complexity $\mathcal{O}((HW)^2)$, which is costly on dense BEV grids ($200 \times 176$). We integrated \textbf{Block-SimAM}~\cite{yang2021simam}, a parameter-free 3D spatial-energy attention module, to provide spatial modulation without learned attention weights.

\paragraph{Design and Placement.}
As illustrated in Figure~\ref{fig:next_s_block}, SimAM is embedded inside each \texttt{MobilePIXORNeXt-\allowbreak S} block across all three stages (C2, C3, C4), for nine attention operators in total. It operates on the intermediate tensor $\mathbf{R} \in \mathbb{R}^{B \times C \times H \times W}$ after projection BN and before LayerScale shortcut addition.

For each channel slice with $M = H \times W$ spatial locations, spatial mean $\mu_c$ and sample variance $\nu_c$ are computed:
\begin{equation}
    \begin{aligned}
        \mu_c &= \frac{1}{M}\sum_{u, v} R_{c, u, v}, \\
        \nu_c &= \frac{1}{\max(M-1, 1)}\sum_{u, v} (R_{c, u, v} - \mu_c)^2.
    \end{aligned}
    \label{eq:spatial_moments}
\end{equation}
Following the energy-minimization formulation of SimAM~\cite{yang2021simam}, the minimized energy $e_t^\star$ and its inverse-energy importance score $s_t$ for neuron $t = R_{c,u,v}$ are:
\begin{equation}
    \begin{aligned}
        e_t^\star &= \frac{4(\nu_c + \lambda_{\mathrm{sim}})}{(t-\mu_c)^2 + 2(\nu_c + \lambda_{\mathrm{sim}})}, \\
        s_t &= \frac{1}{e_t^\star} = \frac{(t - \mu_c)^2}{4(\nu_c + \lambda_{\mathrm{sim}})} + \frac{1}{2}.
    \end{aligned}
    \label{eq:simam_energy}
\end{equation}
Here, $\lambda_{\mathrm{sim}} = 10^{-4}$ is a numerical regularizer. For fixed channel moments, a larger squared deviation $(t-\mu_c)^2$ decreases the minimized energy $e_t^\star$ and increases the inverse-energy importance $s_t$. The sigmoid gate therefore assigns greater weight to more distinctive spatial responses. The attended 3D modulation tensor $\mathbf{A} \in [0, 1]^{B \times C \times H \times W}$ and output $\mathcal{S}(\mathbf{R})$ are:
\begin{equation}
    A_{c, u, v} = \sigma(s_t) = \frac{1}{1 + e^{-s_t}}, \qquad \mathcal{S}(\mathbf{R}) = \mathbf{A} \odot \mathbf{R}.
    \label{eq:simam_modulation}
\end{equation}
Moments and gating scalars are computed in FP32 precision to prevent underflow, then multiplied elementwise with features.

\paragraph{Technical Advantage.}
Block-SimAM adds \textbf{zero parameters} ($0$\,M), whereas stage-level adapters add a learned scalar residual gain $\beta$. Its 3D spatial-and-channel weights modulate sparse LiDAR clusters relative to flat terrain backgrounds. In the ablation, SimAM placement inside each \texttt{MobilePIXORNeXt-S} block improved mAP by $+0.88\,\text{pp}$ over no attention and by up to $+2.10\,\text{pp}$ over stage-level adapters (Table~\ref{tab:attention}).


\subsection{Top-Down Scale-Gated Feature Pyramid Network (SG-FPN Neck)}
\label{sec:neck_architecture}

\paragraph{Motivation.}
Transposed convolutions used for feature upsampling can produce checkerboard artifacts in continuous coordinate regression maps. Bidirectional multiscale networks such as PANet and BiFPN add bottom-up pathways and memory latency. Unweighted addition in conventional FPN can also dilute semantic features. We designed \textbf{Top-Down Scale-Gated Feature Pyramid Network} (\texttt{SG-FPN Neck}) as a unidirectional top-down neck with depthwise scale gating to balance these features at lower cost.

\par\vspace{0.5\baselineskip}
\noindent\begin{minipage}{\columnwidth}
    \centering
    \paperfigure{scale_gated_fusion.pdf}{34mm}{Scale-Gated Feature Pyramid Network (SG-FPN). Top-down features $\mathbf{U}_r$ and lateral features $\mathbf{L}_r$ are dynamically gated via learnable depthwise convolutions to prevent high-level semantic distortion of fine spatial coordinates.}
    \captionof{figure}{Scale-Gated Feature Pyramid Network (SG-FPN) fusion at stride $r \in \{8, 4\}$. Bilinear interpolation upsamples top-down features. Depthwise refinement at stride 8 and projection convolution at stride 4 produce representations $\mathbf{U}_4$ and $\mathbf{U}_3$ without deconvolution checkerboard artifacts. Lateral high-resolution features $\mathbf{L}_r$ are weighted by a zero-initialized depthwise gate $\mathbf{G}_r = 2\sigma(\operatorname{DWConv}_{3\times3}(\mathbf{L}_r + \mathbf{U}_r))$ before element-wise summation $\mathbf{P}_r = \mathbf{U}_r + \mathbf{G}_r \odot \mathbf{L}_r$ to retain boundary detail and suppress background responses.}
    \label{fig:scale_gated_fusion}
\end{minipage}
\par


\paragraph{Mathematical Formulation.}
As shown in Figure~\ref{fig:scale_gated_fusion}, lateral features from Stages C4 ($128$\,ch), C3 ($96$\,ch), and C2 ($48$\,ch) are projected via $1 \times 1$ convolutions into matching lateral dimensions $\mathbf{L}_5 \in \mathbb{R}^{B \times 48 \times 50 \times 44}$, $\mathbf{L}_4 \in \mathbb{R}^{B \times 48 \times 100 \times 88}$, and $\mathbf{L}_3 \in \mathbb{R}^{B \times 32 \times 200 \times 176}$.

Feature aggregation follows a unidirectional top-down path. At stride 8, $2\times$ bilinear interpolation and depthwise $3 \times 3$ refinement with BN and SiLU produce $\mathbf{U}_4 \in \mathbb{R}^{B \times 48 \times 100 \times 88}$:
\begin{align}
    \mathbf{U}_4 &= \operatorname{Refine}_{3\times 3}\left(\operatorname{Bilinear}_{2\times}(\mathbf{L}_5)\right),\nonumber\\
    \mathbf{P}_4 &= \mathbf{U}_4 + \mathbf{G}_4 \odot \mathbf{L}_4,
    \label{eq:scale_gating_u4}
\end{align}
where $\mathbf{G}_4 = 2\sigma(\operatorname{DWConv}_{3\times3}(\mathbf{L}_4 + \mathbf{U}_4))$. At stride 4, $2\times$ bilinear interpolation and standard $3 \times 3$ convolution ($48 \rightarrow 32$) with BN and SiLU produce $\mathbf{U}_3 \in \mathbb{R}^{B \times 32 \times 200 \times 176}$:
\begin{align}
    \mathbf{U}_3 &= \operatorname{Proj}_{3\times 3}\left(\operatorname{Bilinear}_{2\times}(\mathbf{P}_4)\right),\nonumber\\
    \mathbf{P}_3 &= \mathbf{U}_3 + \mathbf{G}_3 \odot \mathbf{L}_3,
    \label{eq:scale_gating_u3}
\end{align}
where $\mathbf{G}_3 = 2\sigma(\operatorname{DWConv}_{3\times3}(\mathbf{L}_3 + \mathbf{U}_3))$. Zero-initialized gating weights and biases make $\mathbf{G}_r$ initially equal to $1.0$, corresponding to standard residual addition. The feature $\mathbf{P}_3$ passes through a $3 \times 3$ convolution, BN, and SiLU to form the shared detector representation $\mathbf{F}_{\mathrm{out}} \in \mathbb{R}^{B \times 32 \times 200 \times 176}$. The neck contains $36,784$ parameters and passes the fused multi-scale features to the prediction towers.

\subsection{Shared Grouped Prediction Towers}
\label{sec:grouped_heads}

\par\vspace{-0.5\baselineskip}
\noindent
\begin{minipage}{\columnwidth}
    \centering
    \paperfigure{geometry_output_2x2.pdf}{75mm}{Geometric parameterization, orientation encoding, and quality-aware score calibration across prediction heads: (a) Continuous subcell coordinate offsets $(\Delta x, \Delta y)$ from anchor point $P$ to object center $C$; (b) Log-transformed bounding box dimension targets $(\log w, \log l)$; (c) Doubled-angle trigonometric yaw representation $(c_2, s_2) = (\cos 2\theta, \sin 2\theta)$ on the unit circle resolving $\pi$-heading ambiguity; and (d) IoU-quality-aware composite score ranking $s = (p_{\mathrm{cls}})^{1-\alpha} q^{\alpha}$ calibrated for post-processing peak selection and rotated NMS.}
    \captionof{figure}{Prediction-head geometry, orientation, and quality-aware scoring. \textbf{(a)} Metric subcell offsets $(\hat{o}_x, \hat{o}_y)$ from grid anchor center $P$ to object center $C$. \textbf{(b)} Log-dimension targets $(t_w, t_l) = (\log w, \log l)$ decoded by exponentiation. \textbf{(c)} Doubled-angle yaw representation $(c_2, s_2) = (\cos 2\theta, \sin 2\theta)$ on the unit circle, resolving $180^\circ$ heading ambiguity through $\pi$-rotational symmetry of BEV footprints. \textbf{(d)} Joint quality scoring combines overlap prediction $\hat{q} \in [0, 1]$ and classification probability $p_{\mathrm{cls}}$ as $s_{\mathrm{det}} = (p_{\mathrm{cls}})^{1-\alpha} \hat{q}^\alpha$ ($\alpha = 0.5$). This score prioritizes well-aligned boxes during peak filtering and rotated NMS.}
    \label{fig:geometry_output}
\end{minipage}
\par

\paragraph{Motivation.}
Road categories have different BEV footprints. Cars are rectangular ($l \approx 4.5$\,m, $w \approx 1.8$\,m), while Pedestrians ($l \approx 0.8$\,m, $w \approx 0.6$\,m) and Cyclists ($l \approx 1.8$\,m, $w \approx 0.7$\,m) are smaller and more deformable. A shared prediction head can introduce gradient competition across these scales; separate heads for every class add parameters. We designed \textbf{Shared Grouped Prediction Towers} (\texttt{GroupedHeader}) to balance class separation and model size.

\paragraph{Topology and Target Formulation.}
We partitioned classes into Group 1 for \texttt{car} ($K_1 = 1$) and Group 2 for \texttt{ped\_cyc} ($K_2 = 2$). For each group $g \in \{1, 2\}$, feature map $\mathbf{F}_{\mathrm{out}}$ passes through an initial $3 \times 3$ Conv-BN-SiLU block with $32 \rightarrow 32$ channels. Five task-specific branches follow:
\begin{itemize}
    \item \texttt{cls}: Classification heatmap $[B, K_g, 200, 176]$ yielding foreground category probability $p_{\mathrm{cls}}$ at grid point $P$ (Fig.~\ref{fig:geometry_output}d).
    \item \texttt{offset}: Metric subcell coordinate offsets $[B,\allowbreak 2,\allowbreak 200,\allowbreak 176]$ predicting $(\hat{o}_x, \hat{o}_y)$, corresponding to continuous displacements $(\Delta x, \Delta y)$ from anchor point $P$ to object center $C$ (Fig.~\ref{fig:geometry_output}a).
    \item \texttt{size}: Bounding box log-dimensions $[B, 2, 200, 176]$ predicting $(\hat{d}_1, \hat{d}_2) = (\log \hat{w}, \log \hat{l})$ (denoted as regression targets $(t_w, t_l)$ in Fig.~\ref{fig:geometry_output}b, decoded via $\exp$).
    \item \texttt{yaw}: Orientation doubled-angle vector $[B, 2, 200, 176]$ predicting $(\hat{a}_1, \hat{a}_2) = (\cos 2\hat{\psi}, \sin 2\hat{\psi})$ (denoted as $(c_2, s_2)$ with heading angle $\theta \equiv \hat{\psi}$ on the unit circle in Fig.~\ref{fig:geometry_output}c to resolve $\pi$-heading ambiguity).
    \item \texttt{iou}: Localization quality score $[B, 1, 200, 176]$ predicting continuous overlap $\hat{q} \in [0, 1]$ ($q(P)$ in Fig.~\ref{fig:geometry_output}d), fused with classification into composite ranking confidence $s_{\mathrm{det}} = (p_{\mathrm{cls}})^{1-\alpha} \hat{q}^\alpha = \sqrt{p_{\mathrm{cls}} \hat{q}}$ ($\alpha=0.5$).
\end{itemize}

For output cell $(u, v)$, the ground-truth geometric targets are formulated as:
\begin{align}
    \mathbf{o}^* &= \begin{pmatrix} x^* - (u s_{\mathrm{out}}\Delta_x + x_{\min}) \\ y^* - (v s_{\mathrm{out}}\Delta_y + y_{\min}) \end{pmatrix},\nonumber\\
    \mathbf{d}^* &= (\log w^*, \log l^*), \qquad \mathbf{a}^* = (\cos 2\psi^*, \sin 2\psi^*).
    \label{eq:gt_targets}
\end{align}
During inference, predicted bounding box parameters are decoded from regression outputs via:
\begin{align}
    \hat{x} &= u s_{\mathrm{out}}\Delta_x + x_{\min} + \hat{o}_x, \quad \hat{y} = v s_{\mathrm{out}}\Delta_y + y_{\min} + \hat{o}_y,\nonumber\\
    \hat{w} &= \exp(\hat{d}_1), \quad \hat{l} = \exp(\hat{d}_2), \quad \hat{\psi} = \frac{1}{2}\operatorname{atan2}(\hat{a}_2, \hat{a}_1).
    \label{eq:decoded_targets}
\end{align}
The doubled-angle formulation $\mathbf{a}^*$ and decoding rule in Eq.~(\ref{eq:decoded_targets}) strictly resolve periodic boundary singularities. Standard scalar regression of heading $\psi \in [-\pi/2, \pi/2)$ suffers from boundary instability: across the periodic boundary $\pm \pi/2$, an infinitesimal geometric rotation yields a jump discontinuity of $\pi$, producing explosive gradient penalties under Smooth-L1 or L2 regression. By projecting $\psi^*$ onto the unit circle $S^1$ via $\mathbf{a}^* = (\cos 2\psi^*, \sin 2\psi^*)$, the mapping $\psi \mapsto \mathbf{a}$ is smooth and inherently periodic with period $\pi$, capturing the geometric $\pi$-rotational symmetry of oriented BEV bounding boxes. At inference, the two-argument arctangent $\operatorname{atan2}(\hat{a}_2, \hat{a}_1)$ is well-defined and continuous everywhere on $\mathbb{R}^2 \setminus \{(0, 0)\}$, ensuring robust angle recovery without branch-cut wrapping artifacts.

Classification supervision utilizes 2D Gaussian heatmaps centered at continuous ground-truth object projections with a minimum radius of 4 cells. Regression targets are assigned within half the Gaussian peak radius, evaluated on decoupled validity masks $\mathbf{M}_{\mathrm{reg}, g}$ for each group $g \in \{1, 2\}$. Specifically, category target ownership is strictly partitioned: anchor cells within the regression radius of a car object activate $\mathbf{M}_{\mathrm{reg}, 1}(u, v) = 1$ to supervise Group 1, while cells within the regression radius of a pedestrian or cyclist activate $\mathbf{M}_{\mathrm{reg}, 2}(u, v) = 1$ to supervise Group 2. Each group tower computes regression losses independently on its respective positive mask, decoupling gradient backpropagation between categories with radically divergent aspect ratios ($4.5$\,m passenger cars vs. $0.8$\,m pedestrians). In the rare event where a vehicle and a vulnerable road user are immediately adjacent such that both masks activate at cell $(u, v)$, both towers regress their respective targets concurrently without inter-head gradient interference. Within Group 2 (\texttt{ped\_cyc}), where pedestrians and cyclists share a single regression head, conflicting assignments on overlapping cells are resolved by assigning ownership to the ground-truth instance with the nearest continuous Euclidean center distance.

\subsection{Multi-Task Geometric Alignment and Quality Supervision (OGA and IQA)}
\label{sec:losses_and_oga}

\paragraph{Motivation.}
We developed \textbf{Oriented Geometric Alignment (OGA)} with \textbf{IoU-Aware Quality Alignment (IQA)} to connect box-parameter regression with oriented overlap. Standard anchor-free regression applies separate Smooth-L1 losses to center $(x, y)$, footprint $(w, l)$, and heading $\psi$. These errors do not vary monotonically with rotated Intersection-over-Union (IoU). Classification confidence can also disagree with localization accuracy, degrading candidate ranking during rotated Non-Maximum Suppression (NMS).

\paragraph{Oriented Geometric Alignment (OGA).}
For each cell $(u, v)$ with $\mathbf{M}_{\mathrm{reg}} = 1$, we decoded predicted and ground-truth boxes into four 2D corners $\{\widehat{\mathbf{c}}_m\}_{m=1}^4$ and $\{\mathbf{c}_m^*\}_{m=1}^4$. Given center $(x, y)$, dimensions $(l, w)$, and longitudinal/lateral unit vectors $\mathbf{n}_{\mathrm{len}} = (\cos\psi, \sin\psi), \mathbf{n}_{\mathrm{wid}} = (-\sin\psi, \cos\psi)$, the four bounding box corners are indexed in quadrant order $[++, +-, -+, --]$:
\begin{equation}
    \begin{aligned}
        \mathbf{c}_1 &= (x, y) + \tfrac{1}{2}l\mathbf{n}_{\mathrm{len}} + \tfrac{1}{2}w\mathbf{n}_{\mathrm{wid}}, \\
        \mathbf{c}_2 &= (x, y) + \tfrac{1}{2}l\mathbf{n}_{\mathrm{len}} - \tfrac{1}{2}w\mathbf{n}_{\mathrm{wid}}, \\
        \mathbf{c}_3 &= (x, y) - \tfrac{1}{2}l\mathbf{n}_{\mathrm{len}} + \tfrac{1}{2}w\mathbf{n}_{\mathrm{wid}}, \\
        \mathbf{c}_4 &= (x, y) - \tfrac{1}{2}l\mathbf{n}_{\mathrm{len}} - \tfrac{1}{2}w\mathbf{n}_{\mathrm{wid}}.
    \end{aligned}
    \label{eq:corner_enumeration}
\end{equation}
OGA combines two geometric objectives:
\par\noindent 1. \textit{Multi-Axis Projection GIoU ($\ell_{\mathrm{proj}}$)}: Projected bounding box overlap is evaluated following Marginalized Generalized IoU~\cite{le2026mgiou} across four unit normal projection axes defining the prediction and ground-truth rectangles:
\begin{equation}
    \mathbf{n}_k \in \left\{ \mathbf{n}_{\mathrm{len}}^{\mathrm{pred}}, \mathbf{n}_{\mathrm{wid}}^{\mathrm{pred}}, \mathbf{n}_{\mathrm{len}}^{\mathrm{gt}}, \mathbf{n}_{\mathrm{wid}}^{\mathrm{gt}} \right\}.
    \label{eq:projection_axes}
\end{equation}
For each axis $\mathbf{n}_k$, 1D intervals of intersection length $I_k$, union length $U_k$, and smallest enclosing hull $E_k$ are determined, yielding 1D GIoU~\cite{rezatofighi2019giou}:
\begin{equation}
    \begin{aligned}
        g_k &= \frac{I_k}{U_k} - \frac{E_k - U_k}{E_k}, \\
        \bar{g} &= \frac{1}{4}\sum_{k=1}^4 g_k, \qquad
        \ell_{\mathrm{proj}} = \frac{1 - \bar{g}}{2}.
    \end{aligned}
    \label{eq:mgiou_loss}
\end{equation}
\par\noindent{\raggedright\textit{2. Scale-Normalized $\pi$-Symmetric Corner Distance ($\ell_{\mathrm{corner}}$)}:\par}\noindent
Under a $\pi$-rotation ($180^\circ$), corner $m$ maps deterministically to corner $5-m$. To ensure accurate vertex localization under doubled-angle symmetry, we evaluate the minimum L1 distance across direct and $\pi$-permuted corner correspondences:
\begin{equation}
    \ell_{\mathrm{corner}} = \frac{\min\left\{
        \begin{array}{l}
            \displaystyle\sum_{m=1}^4 \|\widehat{\mathbf{c}}_m - \mathbf{c}_m^*\|_1, \\
            \displaystyle\sum_{m=1}^4 \|\widehat{\mathbf{c}}_m - \mathbf{c}_{5-m}^*\|_1
        \end{array}
    \right.}{4\max\left(\sqrt{(w^*)^2 + (l^*)^2}, 10^{-6}\right)}.
    \label{eq:corner_loss}
\end{equation}
Normalizing by the target bounding box diagonal ensures scale invariance across categories. The composite geometric loss is $\mathcal{L}_{\mathrm{geo}} = \ell_{\mathrm{proj}} + \beta_{\mathrm{corner}} \ell_{\mathrm{corner}}$ with $\beta_{\mathrm{corner}} = 1.0$.

For positive anchor cells with $\mathbf{M}_{\mathrm{reg}} = 1$ ($N_{\mathrm{pos}} = \sum \mathbf{M}_{\mathrm{reg}}$), the subcell offset, log-size, and doubled-yaw heads are supervised using masked Smooth-$L_1$ loss:
\begin{equation}
    \begin{aligned}
        \mathcal{L}_{\mathrm{task}} &= \frac{1}{N_{\mathrm{pos}} + \epsilon} \sum_{i \in \mathrm{pos}} \operatorname{SmoothL1}(\hat{\mathbf{y}}_i, \mathbf{y}_i^*), \\
        &\quad \text{for } \mathrm{task} \in \{\mathrm{off}, \mathrm{size}, \mathrm{yaw}\},
    \end{aligned}
    \label{eq:smooth_l1_base}
\end{equation}
with threshold $\beta_{\mathrm{reg}} = 1/9$ and $\epsilon = 10^{-6}$.

\paragraph{IoU-Aware Quality Alignment (IQA).}
The quality head predicts continuous overlap score $\hat{q} = \sigma(\hat{Z}_{\mathrm{qual}}) \in [0, 1]$. It is supervised via Binary Cross-Entropy (BCE) over positive anchor cells using the detached MGIoU similarity as the ground-truth target:
\begin{equation}
    \begin{aligned}
        q^* &= \operatorname{stopgrad}(\operatorname{clip}(\bar{g}, 0, 1)), \\
        \mathcal{L}_{\mathrm{qual}} &= \frac{1}{N_{\mathrm{pos}} + \epsilon} \sum_{i \in \mathrm{pos}} \operatorname{BCE}(\hat{q}_i, q_i^*).
    \end{aligned}
    \label{eq:iqa_loss}
\end{equation}
Detaching the target limits gradient instability while retaining supervision for confidence calibration.

\paragraph{Head-Free Alternative: Quality-Coupled OGA (Q-OGA).}
The explicit IQA branch gave the highest ranking accuracy, adding $18{,}626$ parameters ($18.63\,\text{K}$) and $1.25\,\text{ms}$ forward latency. We developed \textbf{Q-OGA} as a head-free alternative for resource-constrained platforms. Q-OGA modulates positive classification targets ($y=1.0$) with detached continuous bounding box IoU quality $q^*$, without an auxiliary head:
\begin{equation}
    y^* = (q^*)^\beta \quad (\beta = 1.0).
    \label{eq:cqfl_target}
\end{equation}
Classification optimization is governed by Quality-Coupled Focal Loss (CQFL):
\begin{align}
    \mathcal{L}_{\mathrm{pos}} &= |y^* - \hat{p}|^\gamma \left[ y^* (-\ln \hat{p}) + (1 - y^*) (-\ln(1 - \hat{p})) \right], \nonumber\\
    \mathcal{L}_{\mathrm{neg}} &= \hat{p}^\gamma (-\ln(1 - \hat{p})),
    \label{eq:cqfl_loss}
\end{align}
where $\gamma = 2.0$. In addition, Q-OGA replaces hard $\pi$-minimum corner distance with a soft-min corner formulation:
\begin{equation}
    \mathcal{D}_{\mathrm{smooth}} = -\tau_{\mathrm{soft}} \ln\left( \mathrm{e}^{-D_{\mathrm{direct}}/\tau_{\mathrm{soft}}} + \mathrm{e}^{-D_\pi/\tau_{\mathrm{soft}}} \right),
    \label{eq:softmin_corner}
\end{equation}
with $\tau_{\mathrm{soft}} = 0.05$, and incorporates Range and Density-Adaptive (RDA) weighting $w(r) = 1.0 + \gamma_{\mathrm{rda}}(r / r_{\max})^{\alpha_{\mathrm{rda}}}$ ($\gamma_{\mathrm{rda}} = 1.5, \alpha_{\mathrm{rda}} = 2.0$). In the evaluation, Q-OGA saved $18{,}626$ parameters and achieved $69.20\%$ Moderate mAP, providing a competitive compact option.

\paragraph{Temperature-Softmax Bounded Uncertainty Weighting (T-SBUW).}
To balance six task losses $\mathcal{T} = \{\mathrm{cls}, \mathrm{off}, \mathrm{size}, \mathrm{yaw}, \mathrm{geo}, \mathrm{qual}\}$ ($J=6$), we developed \textbf{T-SBUW} without manually tuning their relative weights. It uses two coupled mechanisms:
\par\noindent 1. \textit{Bounded Softmax Weights}: Learnable logit parameters $a_k$ (initialized to $0.0$, yielding initial weights $\omega_k = 1.0$) are soft-clamped via hyperbolic tangent with bound $b_{\mathrm{uw}} = 3.0$, and converted to positive normalized weights via temperature-controlled softmax ($\tau_{\mathrm{uw}} = 2.0$):
\begin{equation}
    \begin{aligned}
        \bar{a}_k &= b_{\mathrm{uw}} \tanh(a_k / b_{\mathrm{uw}}), \\
        Z_{\mathrm{uw}} &= \sum_{j \in \mathcal{T}} \exp(\bar{a}_j / \tau_{\mathrm{uw}}), \\
        \omega_k &= J \cdot \frac{\exp(\bar{a}_k / \tau_{\mathrm{uw}})}{Z_{\mathrm{uw}}}, \\
        \sum_{k \in \mathcal{T}} \omega_k &= J.
    \end{aligned}
    \label{eq:tsbuw_weights}
\end{equation}
This construction strictly conserves the total weight.
\par\noindent 2. \textit{Running Scale Calibration}: Exponential Moving Average (EMA) with momentum $\beta_{\mathrm{EMA}} = 0.99$ (initialized with $\mu_k^{(0)} = 1.0$) tracks the unweighted running loss magnitude $\mu_k \leftarrow 0.99\mu_k + 0.01\mathcal{L}_k$ during training. The balanced objective for group $g$ is:
\begin{equation}
    \mathcal{L}_g = \left(\frac{1}{J}\sum_{k \in \mathcal{T}} \mu_k\right) \sum_{k \in \mathcal{T}} \omega_k \frac{\mathcal{L}_k}{\max(\mu_k, 10^{-4})}.
    \label{eq:balanced_loss}
\end{equation}
The total network objective is $\mathcal{L}_{\mathrm{total}} = 0.5\mathcal{L}_1 + 0.5\mathcal{L}_2$. The twelve logit parameters are utilized strictly during training and discarded at inference.


% \newpage
\subsection{Coordinate Decoding and Inference Pipeline}
\label{sec:decoding_and_pipeline}

\noindent
\captionof{algorithm}{End-to-end MobilePIXORNeXt inference and box decoding pipeline.}
\label{alg:inference_pipeline}
\vspace{2pt}
{\small
\noindent\textbf{Require:} Raw LiDAR point cloud $\mathcal{P} = \{(x_i, y_i, z_i, \iota_i)\}_{i=1}^N$, ROI bounds $[x_{\min}, x_{\max}] \times [y_{\min}, y_{\max}] \times [z_{\min}, z_{\max}]$\par
\noindent\textbf{Require:} Score threshold $\tau_{\mathrm{score}} = 0.05$, Rotated NMS threshold $\tau_{\mathrm{nms}} = 0.10$, quality exponent $\alpha = 0.5$\par
\noindent\textbf{Ensure:} Final decoded 3D oriented bounding boxes $\mathcal{B} = \{(c_k, s_k, \hat{x}_k, \hat{y}_k, \hat{l}_k, \hat{w}_k, \hat{\psi}_k)\}_{k=1}^M$\par
\begin{enumerate}
    \setlength{\topsep}{2pt}
    \setlength{\itemsep}{0pt}
    \setlength{\parsep}{0pt}
    \renewcommand{\labelenumi}{\textbf{\arabic{enumi}:}}
    \item Filter raw points $\mathcal{P}$ within inward margin $\epsilon = 0.001$\,m of spatial ROI.
    \item Deterministically rasterize points into 14-channel Strata-BEV14 tensor $\mathbf{X} \in \mathbb{R}^{14 \times 800 \times 704}$ via Eq.~(\ref{eq:density_encoding})--(\ref{eq:offset_encoding}).
    \item Forward $\mathbf{X}$ through ConvNeXt stem to obtain stride-2 feature $\mathbf{F}^{(2)} \in \mathbb{R}^{32 \times 400 \times 352}$.
    \item Extract multi-scale features through Stages 2, 3, 4 via MobilePIXORNeXt-S blocks to obtain $\mathbf{F}^{(4)}, \mathbf{F}^{(8)}, \mathbf{F}^{(16)}$.
    \item Fuse multi-scale features via Top-Down SG-FPN Neck into shared map $\mathbf{F}_{\mathrm{out}} \in \mathbb{R}^{32 \times 200 \times 176}$.
    \item Evaluate Grouped Prediction Towers for Car ($g=1$) and Pedestrian/Cyclist ($g=2$), yielding classification maps $\mathbf{P}_{\mathrm{cls}}$, subcell offsets $\widehat{\mathbf{o}}$, log-sizes $\widehat{\mathbf{d}}$, doubled yaws $\widehat{\mathbf{a}}$, and quality $\widehat{\mathbf{q}}$.
    \item \textbf{for each} class $c \in \{\text{Car}, \text{Pedestrian}, \text{Cyclist}\}$ \textbf{do}
    \item \quad Compute IoU-aware composite ranking map across grid: $\mathbf{S}_{\mathrm{det}, c} = (\mathbf{P}_{\mathrm{cls}, c})^{1 - \alpha} \cdot (\widehat{\mathbf{q}})^\alpha$.
    \item \quad Identify local $3 \times 3$ candidate peaks on composite ranking map $\mathbf{S}_{\mathrm{det}, c}$.
    \item \quad Filter candidate cells satisfying $\mathbf{S}_{\mathrm{det}, c} > \tau_{\mathrm{score}}$ and $\mathbf{P}_{\mathrm{cls}, c} > \tau_{\mathrm{score}}$.
    \item \quad Decode metric coordinates, dimensions, and heading angles via Eq.~(\ref{eq:box_decode}).
    \item \quad Apply Classwise Rotated Non-Maximum Suppression with IoU threshold $\tau_{\mathrm{nms}}$ ranked by $s_{\mathrm{det}}$.
    \item \textbf{end for}
    \item \textbf{return} Final detection set $\mathcal{B}$ sorted in descending order of composite score $s_{\mathrm{det}}$.
\end{enumerate}
}

\par\vspace{0.5\baselineskip}


Algorithm~\ref{alg:inference_pipeline} specifies inference. For candidate cells satisfying the joint score threshold and local peak filtering, bounding boxes are decoded in metric coordinates as:
\begin{align}
    \hat{x} &= x_{\min} + u s_{\mathrm{out}}\Delta_x + \hat{o}_x, \nonumber\\
    \hat{y} &= y_{\min} + v s_{\mathrm{out}}\Delta_y + \hat{o}_y, \nonumber\\
    \hat{w} &= \exp(\hat{d}_1), \qquad \hat{l} = \exp(\hat{d}_2), \nonumber\\
    \hat{\psi} &= \frac{1}{2}\operatorname{atan2}(\hat{a}_2, \hat{a}_1).
    \label{eq:box_decode}
\end{align}
Predicted orientation angle $\hat{\psi} \in [-\pi/2, \pi/2]$ is fully recovered without trigonometric ambiguity.

To align classification confidence with geometric overlap precision, detection candidates are ranked via a joint IoU-aware composite confidence score:
\begin{equation}
    s_{\mathrm{det}} = (p_{\mathrm{cls}})^{1 - \alpha_{\mathrm{score}}} \cdot (\hat{q})^{\alpha_{\mathrm{score}}}, \qquad \alpha_{\mathrm{score}} = 0.5.
    \label{eq:joint_score}
\end{equation}
Candidates with $s_{\mathrm{det}} > 0.05$ undergo class-wise rotated NMS at IoU threshold $\tau_{\mathrm{nms}} = 0.10$. The output is capped at 500 boxes per scan.

\paragraph{Parameter Distribution Breakdown.}
Our complete detector architecture contains \textbf{753,569 deployable parameters} ($< 0.76$\,M, requiring $753{,}581$ trainable parameters during training including loss-balancing logits), with an FP32 checkpoint size of $9.34$\,MB:
\begin{itemize}
    \item \textbf{Backbone Body}: $604,864$ parameters ($80.3\%$).
    \item \textbf{Scale-Gated Neck (SG-FPN)}: $36,784$ parameters ($4.9\%$).
    \item \textbf{Grouped Prediction Heads}: $111,921$ parameters ($14.8\%$), divided into $55,944$ parameters for Group \texttt{car} and $55,977$ parameters for Group \texttt{ped\_cyc}.
    \item \textbf{Training-Only T-SBUW Logits}: 12 parameters ($< 0.01\%$), optimized exclusively during training and eliminated during deployment.
\end{itemize}
The complete pipeline achieved $16.41$\,FPS ($60.96$\,ms) on an NVIDIA Tesla T4 GPU, while in-pipeline model forward execution took $17.88$\,ms ($55.94$\,FPS) and standalone model forward reached $8.67$\,ms ($115.32$\,FPS).
